\documentclass[11pt]{article}
\usepackage[a4paper,margin=1in]{geometry}
\usepackage{amsmath,amssymb,amsthm}
\usepackage{booktabs}
\usepackage{array}
\usepackage{hyperref}

\newtheorem{definition}{Definition}
\newtheorem{theorem}{Theorem}

\newcommand{\XiRMT}{\Xi_{\mathrm{RMT}}}
\newcommand{\CRE}{\mathrm{CRE}}

\title{Step 188: Random Matrix Theory Carrier Pivot}
\author{RH Membrane Adequacy Track}
\date{}

\begin{document}
\maketitle

\section*{Verdict}

\[
\boxed{\mathrm{RMT\_verdict}=V_{\mathrm{RMT\_outside\_dichotomy\_scope}}.}
\]

Random Matrix Theory is a statistical evidence and consistency-checking
carrier for zeta-zero data.  It is not a closure carrier for RH and therefore
does not refute the Framework RH-Carrier Dichotomy.  Instead it refines the
domain: the dichotomy applies to RH-analogous closure carriers, not to
statistical evidence carriers.

\section*{Inherited Record Audit}

The inherited records do not provide an accepted RMT carrier package.  The
closest local records are the Hilbert--Polya/Berry--Keating pivot and the
Step 187 dichotomy theorem.  Those records concern spectral closure carriers,
whereas RMT supplies statistical laws for spectra.  Therefore the RMT carrier
must be declared directly from the classical RMT literature.

\section*{Carrier Declaration}

\begin{definition}[RMT statistical comparison carrier]
Let
\[
C_{\mathrm{RMT}}(N,T)=(M_N,Z_T),
\]
where \(M_N\) is either a GUE Hermitian random matrix or a CUE Haar-random
unitary matrix, and
\[
Z_T=\{\gamma_j:0<\gamma_j\le T\}
\]
is the zeta-zero ordinate data after the critical-line normalization has been
chosen.  The matrix side is scaled to mean spacing one, and the zeta side is
locally normalized by the usual zero-density factor.
\end{definition}

The native probes are:
\begin{itemize}
\item pair correlation \(R_2\);
\item \(n\)-level correlations \(R_n\);
\item nearest-neighbor spacing statistics;
\item characteristic-polynomial moment statistics.
\end{itemize}
The dissolving probes are computed zero ordinates and empirical
spacing/correlation/moment data.

\section*{RMT Residual}

Define \(\XiRMT\) as a statistical distance between RMT predictions and
zeta-zero data.  Examples include:
\[
\|\rho^{\zeta}_2-\rho^{\mathrm{GUE}}_2\|_{L^2(w)}^2,
\]
a Kolmogorov--Smirnov distance for spacing distributions, \(n\)-level
correlation discrepancy, or moment discrepancy between zeta moments and
characteristic-polynomial moment predictions.

Thus
\[
\XiRMT=0
\]
means exact or asymptotic statistical agreement with the chosen random-matrix
ensemble.  It does not mean that RH has been proved.

\section*{CRE Audit}

\begin{center}
\begin{tabular}{p{0.30\linewidth}p{0.18\linewidth}p{0.42\linewidth}}
\toprule
Question & Answer & Reason\\
\midrule
Is \(\XiRMT=0\) equivalent to RH? & no
& exact GUE/CUE agreement is statistical agreement, not a theorem that all
zeros lie on the critical line\\
Does RMT require RH to state cleanly? & yes
& standard normalized zero statistics use an ordered critical-line sequence
\(\gamma_j\)\\
Is RMT a non-CRE proved RH analog? & no
& it is not the zeta function of a separate object with a proved RH analog\\
CRE status & not closure carrier
& RMT is evidence/consistency machinery, not a residual-closing carrier\\
\bottomrule
\end{tabular}
\end{center}

Montgomery's pair-correlation theorem/conjecture is not used as an
unconditional RH proof.  In the standard form relevant here it is conditional
on RH and restricted test-function support.  Odlyzko's computations are
powerful numerical evidence, and Keating--Snaith supplies moment predictions,
but neither is a closure theorem for RH.

\section*{Dichotomy Applicability}

The Step 187 dichotomy classifies typed carriers aimed at closing an
RH-analogous statement about an \(L\)-function or zeta function.  RMT does not
define such an \(L_C\) with a native RH statement.  It defines statistical
comparators for Riemann-zero data.

Therefore
\[
\boxed{\text{RMT is outside the dichotomy scope, not a refutation.}}
\]
The coverage conjecture remains unchanged on closure carriers, with the
domain refined to exclude statistical evidence carriers.

\section*{Classical References}

\begin{itemize}
\item H. Montgomery, 1973, pair correlation of zeta zeros: pair-correlation
theorem/conjecture in RH-conditional restricted-test context.
\item A. Odlyzko, numerical computations of high zeta zeros and GUE
statistics.
\item J. P. Keating and N. C. Snaith, 2000: CUE characteristic-polynomial
model for zeta moments.
\item P. Diaconis and M. Shahshahani, 1994: random unitary matrix trace and
eigenvalue statistics.
\end{itemize}

\section*{Nonclaim Boundary}

This step does not prove RH and does not claim that RMT proves or disproves
RH.  It does not treat pair correlation as an unconditional proof.  It does
not claim exact GUE/CUE agreement implies RH without a separate critical-line
theorem.  It only classifies RMT as a statistical evidence carrier outside the
strict closure-carrier domain of the dichotomy.

\end{document}
